% ============ 第1章 §1.7 Hilbert变换的L^p有界性(原书 pp.51-54 = PDF pp.63-66) ============
\section{\texorpdfstring{Hilbert 变换的 $L^p$ 有界性($1<p<\infty$)}{Hilbert 变换的 Lp 有界性(1<p<∞)}}\label{sec:1.7}
\index{Hilbert 变换!Lp 有界性}


关于 Hilbert 变换的 $L^{p}$ 有界性,$L^{2}$ 有界性已经证明,下面先处理 $2 < p < \infty$ 的情形.

\begin{proposition}\label{prop:1.7.1}
设 $2 \leq p < \infty$.存在常数 $C_{p}$ 使得
\begin{equation*}
\| Hf \|_{L^{p}} \leq C_{p}\| f \|_{L^{p}} \quad (f \in \mathcal{S}(\mathbb{R})).
\end{equation*}
\end{proposition}

\begin{proof}
由于 $f \in \mathcal{S}(\mathbb{R})$,由引理~\ref{lem:1.3.2},$|Hf(x)| \leq C(1 + |x|)^{-1}$.因此,容易看出
\begin{equation*}
M(Hf)(x) \leq \frac{C}{1 + |x|}\log(1 + |x|) \in L^{r}(\mathbb{R}) \quad (1 < r < \infty).
\end{equation*}
由此,可用定理~\ref{thm:1.6.8},
\begin{equation*}
\| M(Hf) \|_{L^{p}} \leq C_{p}\| (Hf)^{\#} \|_{L^{p}} \quad (1 < p < \infty).
\end{equation*}
结合命题~\ref{prop:1.4.5},
\begin{equation*}
\| M(Hf) \|_{L^{p}} \leq C_{p}\| M_{2}f \|_{L^{p}}.
\end{equation*}
进一步使用推论~\ref{cor:1.5.6},当 $p > 2$ 时,
\begin{equation*}
\| Hf \|_{L^{p}} \leq C_{p}'\| f \|_{L^{p}}.
\end{equation*}
\end{proof}

由 $2 \leq p < \infty$ 时的 $L^{p}$ 有界性与对偶性,可证明 $1 < p < \infty$ 时的 $L^{p}$ 有界性.

\begin{theorem}\label{thm:1.7.2}
设 $1 < p < \infty$.存在常数 $C_{p}$ 使得
\begin{equation*}
\| Hf \|_{L^{p}} \leq C_{p}\| f \|_{L^{p}} \quad (f \in \mathcal{S}(\mathbb{R})).
\end{equation*}
\end{theorem}

\begin{proof}
$2 \leq p < \infty$ 的情形已经证明,故设 $1 < p < 2$.对 $f, g \in \mathcal{S}(\mathbb{R})$,
\begin{align*}
(Hf, g) & = \int_{-\infty}^{\infty} Hf(x)\overline{g(x)}\dd x \\
& = \frac{1}{2\pi}\int_{-\infty}^{\infty}\widehat{Hf}(\xi)\overline{\hat{g}(\xi)}\dd \xi \qquad \text{(Parseval)} \\
& = \frac{1}{2\pi}\int_{-\infty}^{\infty} -i\frac{\xi}{|\xi|}\hat{f}(\xi)\overline{\hat{g}(\xi)}\dd \xi \\
& = \frac{1}{2\pi}\int_{-\infty}^{\infty} \hat{f}(\xi)\overline{i\frac{\xi}{|\xi|}\hat{g}(\xi)}\dd \xi \\
& = \int_{-\infty}^{\infty} f(x)\overline{(-1)Hg(x)}\dd x = (f, -Hg) \qquad \text{(Parseval)}
\end{align*}
设 $p'$ 是满足 $1/p + 1/p' = 1$ 的数,则 $2 < p' < \infty$,故使用命题~\ref{prop:1.7.1},
\begin{equation*}
|(Hf, g)| = |(f, -Hg)| \leq \| f \|_{L^{p}}\| Hg \|_{L^{p'}} \leq C_{p'}\| f \|_{L^{p}}\| g \|_{L^{p'}}.
\end{equation*}
因此,由引理~\ref{lem:1.3.2} 有 $Hf \in L^{p}(\mathbb{R})$,且由于 $\mathcal{S}(\mathbb{R})$ 在 $L^{p'}(\mathbb{R})$ 中稠密,
\begin{equation*}
\left|\int_{-\infty}^{\infty} Hf(x)\overline{g(x)}\dd x\right| \leq C_{p'}\| f \|_{L^{p}}\| g \|_{L^{p'}} \quad (g \in L^{p'}(\mathbb{R})).
\end{equation*}
因此,由 $L^{p}$ 与 $L^{p'}$ 的对偶性,
\begin{equation*}
\| Hf \|_{L^{p}} \leq C_{p'}\| f \|_{L^{p}}.
\end{equation*}

(不使用对偶性,注意到由引理~\ref{lem:1.3.2} 有 $|Hf(x)| \leq C(1 + |x|)^{-1}$,在最后一个不等式中令 $g(x) = |Hf(x)|^{p-1}\frac{Hf(x)}{|Hf(x)|}\chi_{\operatorname{supp}(Hf)}(x)$,也可以直接验证 $\| Hf \|_{L^{p}} \leq C_{p'}\| f \|_{L^{p}}$.)
\end{proof}

\begin{note}\label{note:1.7.3}
由上面的定理,由于 $\mathcal{S}(\mathbb{R})$ 在 $L^{p}(\mathbb{R})$ 中稠密,Hilbert 变换可唯一地扩张到整个 $L^{p}(\mathbb{R})$;把扩张后的变换仍记作 $Hf$,则 Hilbert 变换是从 $L^{p}(\mathbb{R})$ 到 $L^{p}(\mathbb{R})$ 的有界线性算子,且
\begin{equation*}
\| Hf \|_{L^{p}} \leq C_{p}\| f \|_{L^{p}} \quad (f \in L^{p}(\mathbb{R})).
\end{equation*}
\end{note}

\begin{corollary}\label{cor:1.7.4}
设 $1 < p < \infty$.则存在常数 $C_{p}$ 使得
\begin{equation*}
(Hf)^{\#}(x) \leq C_{p}M_{p}f(x) \quad (x \in \mathbb{R},\ f \in \mathcal{S}(\mathbb{R})).
\end{equation*}
\end{corollary}

\begin{proof}
对命题~\ref{prop:1.4.5} 的证明稍加修改即可证明.事实上,在命题~\ref{prop:1.4.5} 的证明中,
\begin{equation*}
I_{2} \leq CM_{1}f(x).
\end{equation*}
关于 $I_{1}$,
\begin{align*}
& I_{1} \leq \left(\frac{1}{|Q|}\int_{Q}|Hf_{1}(x)|^{p}\dd x\right)^{1/p} \quad \text{(Hölder)} \\
& \quad \leq \left(\frac{1}{|Q|}\int_{-\infty}^{\infty}|Hf_{1}(x)|^{p}\dd x\right)^{1/p} \\
& \quad \leq C_{p}\left(\frac{1}{|Q|}\int_{-\infty}^{\infty}|f_{1}(x)|^{p}\dd x\right)^{1/p} \quad \text{(Hilbert 变换的 $L^{p}$ 有界性)} \\
& \quad = C_{p}\left(\frac{|Q'|}{|Q|}\cdot\frac{1}{|Q'|}\int_{Q'}|f(x)|^{p}\dd x\right)^{1/p} \\
& \quad \leq 2^{1/p}C_{p}M_{p}f(x_{0}).
\end{align*}
其余与命题~\ref{prop:1.4.5} 的证明相同.
\end{proof}
